% !TEX program = pdflatex
\documentclass[12pt]{article}
\usepackage[a4paper,margin=1in]{geometry}
\usepackage[T1]{fontenc}
\usepackage[utf8]{inputenc}
\usepackage{lmodern}
\usepackage{amsmath,amssymb,amsthm,mathtools,bm}
\usepackage{booktabs,array,tabularx}
\usepackage{enumitem}
\newcolumntype{L}[1]{>{\raggedright\arraybackslash}p{#1}}
\newcolumntype{Y}{>{\raggedright\arraybackslash}X}
\usepackage{setspace}
\usepackage[round,authoryear]{natbib}
\usepackage[colorlinks=true,linkcolor=blue,citecolor=blue,urlcolor=blue]{hyperref}

\onehalfspacing
\setlength{\parindent}{1.5em}
\setlength{\parskip}{0.12em}
\setlist[itemize]{leftmargin=2em,itemsep=0.25em,topsep=0.35em}
\setlist[enumerate]{leftmargin=2em,itemsep=0.25em,topsep=0.35em}
\emergencystretch=3em

\newtheorem{proposition}{Proposition}
\newtheorem{corollary}{Corollary}
\newtheorem{remark}{Remark}
\newtheorem{definition}{Definition}
\newtheorem{assumption}{Assumption}

\newcommand{\E}{\mathbb{E}}
\newcommand{\1}{\mathbf{1}}
\newcommand{\dd}{\mathrm{d}}
\newcommand{\calR}{\mathcal{R}}

\title{Regulatory Separation and Pharmaceutical Innovation: Evidence from China's MAH Reform}
\author{}
\date{}

\begin{document}
\maketitle
\thispagestyle{plain}

\begin{abstract}
\noindent
This paper develops a partial-equilibrium mechanism model of how China's
Marketing Authorization Holder (MAH) system changes the value of advancing
viable drug projects and their commercialization organization. Developers
choose one common original-drug innovation investment / project-advancement
intensity before project characteristics and organizational routes are
realized. MAH makes retained entrusted manufacturing legally available while
the developer retains authorization-holder responsibility. Developers also
differ in predetermined financing capacity for commercialization commitments.
Internal production and retained entrusted production require distinct
up-front commitments, but those commitments are liquidity requirements rather
than additional route costs. MAH does not expand financing supply. It reduces
the amount of vertical-integration capital that must be financed to retain
commercialization rights, and is therefore especially valuable when a
developer cannot finance full internalization but can finance retained
entrusted production. Route choice remains deterministic, and the price of
qualified manufacturing services is endogenous. Reform effects are confined
to project--developer pairs for which the new route is financeable and improves
on the best old option; service scarcity can attenuate these gains. The
baseline does not treat MAH as a demand, credit-supply,
scientific-productivity, downstream-success, or patenting shock.
\end{abstract}

\noindent\textit{Keywords:} MAH; pharmaceutical innovation; commercialization frictions; complementary assets; entrusted production; financing capacity.\\
\noindent\textit{JEL codes:} D22; L65; O31; O38; I18.

\section{Introduction}

China's MAH system separates the right to hold marketing authorization from a
strict requirement of in-house production. This institutional change is
potentially important for drug development because many research-originating
firms can generate valuable projects but lack complete manufacturing
capacity. The central mechanism studied here is therefore not demand expansion
or scientific productivity. It is an expected-commercialization-value
channel: when a firm anticipates that an advancing project can later use
qualified external manufacturing while the firm retains authorization, the
expected value of advancing that project may increase. The reform can be
particularly important when weak internal manufacturing capability makes
vertical integration capital intensive. MAH does not relax the firm's
financing constraint; it changes the organizational technology required to
retain commercialization rights.

The mechanism builds on two established ideas. First, pharmaceutical
innovation responds to expected downstream rewards
\citep{Schmookler1966,AcemogluLinn2004,Finkelstein2004,DuboisDeMouzonScottMortonSeabright2015,BudishRoinWilliams2015}.
Second, an innovator's ability to appropriate returns depends on complementary
assets, firm boundaries, and access to specialized inputs
\citep{Teece1986,AroraFosfuriGambardella2001,GansStern2003,GrossmanHelpman2002,BolerMoxnesUlltveitMoe2015,Ma2026}.
Financing frictions can make large up-front innovative investments difficult
to fund \citep{HallLerner2010}. MAH combines these ideas narrowly: it changes
how an innovator can access manufacturing capability while retaining the
authorization asset, without changing the supply of finance.

Recent evidence on China's pharmaceutical reforms helps separate this route
mechanism from two nearby channels. \citet{JiaMaYangZhang2023} study the 2015
drug-review reform and emphasize approval delay, review capacity, and the
composition of firms and products entering the pipeline.
\citet{BarwickXiaXia2026} study the National Reimbursement Drug List and
emphasize demand and market-size incentives. Those channels can coexist with
MAH, but they are not the MAH mechanism in this paper. The baseline holds
scientific productivity and the market-return environment fixed and asks
whether a regulated holder--producer separation route changes expected
commercialization value.

Closely related evidence also shows why project advancement and upstream
patenting must remain distinct. \citet{Gu2024} finds that allowing production
outsourcing raises clinical-trial registrations among innovators without
production facilities, while patent applications fall among firms with fewer
financial resources. The paper also reports heterogeneous responses across
development stages and drug classes. Gu studies financing-induced reallocation
between research and development. The present baseline instead allows
financing capacity to determine whether a developer can fund vertical
integration or the smaller organizational commitment required for retained
entrusted production. This distinction is a maintained model boundary, not a
claim of empirical novelty: MAH changes route feasibility and value here, but
does not directly constrain the advancement control or generate patents.

For a development-economics audience, the relevant bottleneck is a regulated
complementary asset. MAH can relax a manufacturing bottleneck for
research-oriented firms while preserving holder responsibility. Scarcity in
qualified production support can still attenuate the effect.

The model is deliberately narrow. It is not a full dynamic structural model
of the pharmaceutical industry. It does not solve product-market clearing,
an input market for project advancement, free entry, exit, or an invariant
firm distribution. The baseline closes only the market directly tied to the
MAH mechanism: qualified manufacturing services. Reform-induced demand for
retained entrusted production can raise the equilibrium service price
\(p_m^*\), so the model does not mechanically imply a positive effect for
every project or developer.

The contribution is to organize the mechanism in a way that is internally
consistent with feasible data. The theory separates project advancement from
planning-stage project arrival, deterministic organizational assignment,
observed holder--producer separation, and realized products. The empirical
role of the model is correspondingly modest: it specifies empirically interpretable
interfaces for project-level organization, retained outcomes, and service
scarcity without claiming that approval-side data separately identify every
primitive friction.

The main text gives the shortest complete commercialization baseline. The
technical appendix records primitives, financing and payoff accounting,
pricing derivations, deterministic sorting, project advancement,
qualified-capacity equilibrium, comparative statics, observed outcomes, and
explicitly separated inactive extensions.

\section{Institutional Background}

The relevant institutional point is timing. Under China's drug-registration
framework, a firm need not already be the final commercial holder when it
begins to advance a viable project. The economic choice is forward-looking: a
firm invests today while anticipating the manufacturing and commercialization
environment it will face if a project advances. Production support and
manufacturing readiness must be planned before final approval is observed.
After obtaining a Drug Approval License, the applicant becomes the MAH
\citep{NMPARegistration2022}. The model therefore treats MAH as changing the
anticipated availability and value of a future organizational route, not
current scientific productivity or downstream success.

The MAH reform creates a regulated holder-producer separation route. The 2016
pilot allowed eligible applicants in selected regions to become holders under
the pilot arrangement
\citep{StateCouncilMAHPilot2016,CFDAMAHImplementation2016}. The later legal
framework consolidated the principle that an MAH may manufacture by itself or
entrust production to a qualified manufacturer
\citep{DrugAdministrationLaw2019}. Subsequent NMPA rules require holder-side
quality responsibility, assessment of the entrusted manufacturer, quality and
entrustment agreements, and ongoing supervision
\citep{NMPAHolderResponsibility2022,NMPAManufacturing2022,NMPAEntrustedProduction2023}.

These rules imply two modeling restrictions. First, entrusted production is
not anonymous outsourcing or transfer. It is a regulated retained route in
which authorization remains with the developer while production is performed
by a qualified external manufacturer. Second, MAH does not eliminate holder
responsibility. The entrusted payoff therefore retains a holder-side burden
\(\mu_E\), while the institutional wedge \(\tau_E(M)\) governs effective
route availability. The downstream realization probability \(s(q)\) is
route independent and invariant to MAH in the baseline.

\section{Model}
\label{sec:p16m-model}

\subsection{Environment and timing}
\label{subsec:p16m-environment}

There is a continuum of developers. Developer \(i\) has predetermined
project-advancement capability \(a_i>0\), internal manufacturing capability
\(k_i>0\), and commercialization-stage financing capacity \(\ell_i>0\):
\begin{equation}
  \theta_i=(a_i,k_i,\ell_i)\sim H(a,k,\ell).
  \label{eq:p16m-developer-type}
\end{equation}
The joint distribution permits arbitrary correlation. Financing capacity is
predetermined and measures the largest up-front organizational commitment the
developer can fund; it is distinct from both research capability and
manufacturing know-how. A viable project that reaches
commercialization-relevant planning draws value \(q>0\) and manufacturing
requirement \(m>0\) from the exogenous distribution \(F(q,m)\). For empirical
decompositions only,
\[
  F(q,m)=\sum_{g\in\{O,\mathrm{Inc}\}}\rho_gF_g(q,m),
  \qquad
  \rho_g\geq0,\quad \sum_g\rho_g=1.
\]
The class label \(g\) creates no class-specific control.

The institutional regime is \(M\in\{0,1\}\). The pre-MAH regime sets the
barrier to retained entrusted manufacturing at infinity; the MAH regime
makes it finite:
\begin{equation}
  \tau_E(0)=+\infty,
  \qquad
  \tau_E(1)=\bar\tau_E<+\infty,
  \qquad \bar\tau_E\geq0.
  \label{eq:p16m-policy}
\end{equation}
The available routes are internal manufacturing \(I\), retained entrusted
manufacturing \(E\), non-retained transfer \(T\), and abandonment or delay
\(A\). Under \(E\), the developer remains the authorization holder and a
qualified external producer manufactures. Thus \(E\) is not transfer.

Before project characteristics and routes are realized, developer \(i\)
chooses one common \(x_i\geq0\), interpreted as original-drug innovation
investment / project-advancement intensity. It is broader than pure clinical
development, but excludes basic research, compound discovery,
patent-generating effort, and patent applications. The expected measure of
viable projects reaching route planning is
\begin{equation}
  \lambda_i^{\mathrm{plan}}=a_ix_i.
  \label{eq:p16m-planning}
\end{equation}

Timing is as follows. First, \(M\) is observed and developers anticipate the
market-consistent qualified-service price. Second, after observing
\((a_i,k_i,\ell_i,M)\), they choose \(x_i\). Third, planning-stage projects
draw \((q,m)\); the developer checks route financing requirements against
\(\ell_i\) and chooses among the financeable routes. Fourth, downstream
realization occurs with exogenous probability \(s(q)\). Finally, aggregate
capacity demand and supply must validate the anticipated service price. The
last step is an equilibrium consistency condition, not a late unanticipated
shock. There is no financing decision or finance market. Upstream research
and patent generation may help determine predetermined \(a_i\), but they are
outside the baseline decision problem.

\subsection{Commercialization technologies}
\label{subsec:p16m-technologies}

Conditional on successful commercialization, residual demand is
\(y(p;q)=Aqp^{-\varepsilon}\), where \(A>0\) and \(\varepsilon>1\).
Product-price optimization gives
\[
  p^*(c)=\frac{\varepsilon}{\varepsilon-1}c
\]
and the gross operating present value
\begin{equation}
  R(q,c)
  =
  \frac{Aq}{1-\beta\varphi}
  \frac{(\varepsilon-1)^{\varepsilon-1}}
       {\varepsilon^\varepsilon}
  c^{1-\varepsilon},
  \qquad
  \beta\in(0,1),\quad \varphi\in[0,1).
  \label{eq:p16m-return}
\end{equation}
The appendix derives the FOC, SOC, and accounting. In particular,
\(R_q>0\) and \(R_c<0\).

Internal production uses marginal cost \(c_I(m,k_i)\), with
\(c_{I,m}>0\) and \(c_{I,k}<0\), and setup cost \(F_I(m,k_i)\), with
\(F_{I,m}>0\) and \(F_{I,k}<0\). Both are finite for \(k_i>0\). Its economic
route value is
\begin{equation}
  W_i^I(q,m)
  =
  s(q)R\!\left(q,c_I(m,k_i)\right)-F_I(m,k_i).
  \label{eq:p16m-internal}
\end{equation}
Internalization requires an up-front financing commitment \(J_I(m,k_i)>0\),
where
\[
  J_{I,m}>0,
  \qquad
  J_{I,k}<0.
\]
The commitment is a liquidity requirement, not an additional real cost:
\(F_I\) is subtracted once in \(W_i^I\), whereas \(J_I\) only determines
whether that route can be funded.

Entrusted production uses external marginal cost \(c_E(m)\), requires
\(b(m)>0\) units of qualified capacity with \(b'(m)>0\), incurs setup cost
\(F_E(m)\), and leaves the holder with burden \(\mu_E\geq0\). At the
qualified-capacity price \(p_m\),
\begin{equation}
  \begin{aligned}
  W_i^E(q,m;M,p_m)
  ={}&
  s(q)R\!\left(q,c_E(m)\right)-F_E(m)-p_mb(m)\\
  &-\mu_E-\tau_E(M).
  \end{aligned}
  \label{eq:p16m-entrusted}
\end{equation}
Retained entrusted production requires its own positive up-front commitment
\(J_E(m)\), with \(J_E'(m)>0\). This is also a liquidity requirement and is
not subtracted from \(W_i^E\). Each real monetary cost therefore enters once.
The transfer and abandonment values are
\begin{equation}
  W^T(q,m)=T(q,m),
  \qquad
  W^A=0.
  \label{eq:p16m-outside}
\end{equation}

Financing-adjusted route values are
\begin{equation}
  \widetilde W_i^I(q,m;\ell_i)
  =
  \begin{cases}
    W_i^I(q,m), & J_I(m,k_i)\leq\ell_i,\\
    -\infty, & J_I(m,k_i)>\ell_i,
  \end{cases}
  \label{eq:p16m-finance-adjusted-internal}
\end{equation}
and
\begin{equation}
  \widetilde W_i^E(q,m;M,p_m,\ell_i)
  =
  \begin{cases}
    W_i^E(q,m;M,p_m), & J_E(m)\leq\ell_i,\\
    -\infty, & J_E(m)>\ell_i.
  \end{cases}
  \label{eq:p16m-finance-adjusted-entrusted}
\end{equation}
MAH changes none of \(\ell_i,J_I,J_E\). Pre-MAH unavailability of \(E\)
continues to arise from \(\tau_E(0)=+\infty\), not from financing.

\subsection{Organization and project advancement}
\label{subsec:p16m-organization}

At a fixed \(p_m\), the project chooses deterministically:
\begin{equation}
  \begin{aligned}
  \widetilde W_i(q,m;M,p_m,\ell_i)
  &=
  \max\{\widetilde W_i^I,\widetilde W_i^E,W^T,W^A\},\\
  r_i^*(q,m;M,p_m,\ell_i)
  &\in
  \arg\max_{r\in\{I,E,T,A\}} \widetilde W_i^r.
  \end{aligned}
  \label{eq:p16m-route-choice}
\end{equation}
Continuous heterogeneity makes ties a zero-measure event; no logit or
inclusive value is used.

The internal-minus-entrusted gap is
\begin{equation}
  \begin{aligned}
  \Delta_{IE}(k_i)
  ={}&
  s(q)\!\left[
  R\!\left(q,c_I(m,k_i)\right)-R\!\left(q,c_E(m)\right)
  \right]-F_I(m,k_i)\\
  &+F_E(m)+p_mb(m)+\mu_E+\tau_E(M).
  \end{aligned}
  \label{eq:p16m-gap}
\end{equation}
Holding \((q,m,M,p_m)\) fixed on the internally feasible domain,
\[
  \Delta_{IE,k}
  =
  s(q)R_c\!\left(q,c_I\right)c_{I,k}-F_{I,k}>0.
\]
\begin{proposition}[Organizational feasibility and sorting]
\label{prop:main-finance-sorting}
Route \(I\) is financially feasible if and only if
\(J_I(m,k_i)\leq\ell_i\); route \(E\) is financially feasible if and only if
\(J_E(m)\leq\ell_i\), in addition to its institutional availability. Since
\(J_{I,k}<0\), stronger internal capability weakly expands the set of projects
for which internalization can be financed. Conditional on both retained
routes being financeable, on both beating \(T\) and \(A\), and on the endpoint
crossing conditions, there is a unique value cutoff
\(k^*(q,m;p_m,M)\) satisfying \(\Delta_{IE}(k^*)=0\). Developers below it
select \(E\), and those above it select \(I\).
\end{proposition}

The financeability screen precedes value ranking; comparative-advantage
sorting remains. In the empirically relevant low-capability region, impose
the local condition
\[
  J_E(m)<J_I(m,k_i).
\]
It says that retained outsourcing requires less up-front financing than
building internal capacity and is not a global restriction.

Define the best pre-MAH value and post-MAH value by
\begin{equation}
  \begin{aligned}
  W_i^0(q,m;\ell_i)
  &=\max\{\widetilde W_i^I(q,m;\ell_i),W^T(q,m),0\},\\
  W_i^1(q,m;p_m,\ell_i)
  &=\max\{W_i^0(q,m;\ell_i),
           \widetilde W_i^E(q,m;1,p_m,\ell_i)\}.
  \end{aligned}
  \label{eq:p16m-old-new-values}
\end{equation}
The financing-adjusted MAH-relevant set is
\begin{equation}
  \mathcal C_i(p_m,\ell_i)
  =\left\{(q,m):
  J_E(m)\leq\ell_i,\quad
  W_i^E(q,m;1,p_m)>W_i^0(q,m;\ell_i)
  \right\}.
  \label{eq:p16m-relevant-set}
\end{equation}

\begin{proposition}[MAH-relevant projects and the financing corridor]
\label{prop:main-financing-corridor}
At a fixed \(p_m\),
\begin{equation}
  W_i^1-W_i^0
  =
  [\widetilde W_i^E-W_i^0]_+,
  \label{eq:p16m-positive-part}
\end{equation}
so a strict project-value gain occurs exactly on
\(\mathcal C_i(p_m,\ell_i)\). Suppose locally that
\(J_E(m)<J_I(m,k_i)\). If \(\ell_i<J_E(m)\), neither retained route is
financeable and MAH has no effect. If
\(J_E(m)\leq\ell_i<J_I(m,k_i)\), only retained entrusted production is
financeable, and the gain is strict exactly when
\(W_i^E>\max\{W^T,0\}\). If \(\ell_i\geq J_I(m,k_i)\), both retained routes
are financeable and the gain is
\([W_i^E-\max\{W_i^I,W^T,0\}]_+\). The reform effect is therefore not
globally monotone in financing capacity.
\end{proposition}

Expected optimized value per planning-stage project is
\begin{equation}
  \Omega_i(M,p_m)
  =
  \int \widetilde W_i(q,m;M,p_m,\ell_i)\,dF(q,m).
  \label{eq:p16m-omega}
\end{equation}
Developer \(i\) solves
\begin{equation}
  \max_{x_i\geq0}
  \left\{
  \beta a_ix_i\Omega_i(M,p_m)
  -
  \frac{\kappa}{1+\nu}x_i^{1+\nu}
  \right\},
  \qquad \kappa>0,\quad \nu>0.
  \label{eq:p16m-advancement-problem}
\end{equation}
Strict concavity gives the unique solution
\begin{equation}
  x_i^*(M,p_m)
  =
  \left[
  \frac{\beta a_i}{\kappa}\Omega_i(M,p_m)
  \right]^{1/\nu},
  \label{eq:p16m-advancement}
\end{equation}
including \(x_i^*=0\) when \(\Omega_i=0\). The binary reform affects
\(x_i^*\) only through \(\Omega_i\). At a fixed \(p_m\), its effect is
strict only if \(E\) improves on the pre-MAH maximum on a
positive-probability project set.

\begin{proposition}[Project advancement and heterogeneity]
\label{prop:main-finance-advancement}
At a fixed \(p_m\), the reform strictly raises \(x_i^*\) if and only if
\begin{equation}
  \int_{\mathcal C_i(p_m,\ell_i)}
  [W_i^E(q,m;1,p_m)-W_i^0(q,m;\ell_i)]\,dF(q,m)>0.
  \label{eq:p16m-strict-advancement}
\end{equation}
Research capability \(a_i\) scales any positive response. Since
\(J_{I,k}<0\), lower \(k_i\) makes the financing corridor more likely at a
given \(\ell_i\), but it does not by itself guarantee a treatment effect.
Within either regime, a larger \(\ell_i\) weakly expands the feasible route
set and hence weakly raises \(x_i^*\). The reform response
\(\Delta x_i(\ell_i)\), however, has no global monotonic sign because greater
financing capacity also restores internalization as a pre-MAH option.
\end{proposition}

\subsection{Qualified manufacturing-service equilibrium}
\label{subsec:p16m-cmo}

Qualified supplier \(j\), with efficiency \(z_j\), chooses capacity
\(s_j\geq0\):
\begin{equation}
  \max_{s_j\geq0}\{p_ms_j-\Psi(s_j;z_j)\}.
  \label{eq:p16m-supplier}
\end{equation}
With \(\Psi_{ss}>0\), positive-price capacity is uniquely determined by
\(p_m=\Psi_s(s_j^*;z_j)\), and aggregate supply is
\[
  S_m(p_m)=\int s_j^*(p_m,z)\,dH_C(z).
\]

Expected entrusted capacity per planning-stage project is
\begin{equation}
  \chi_i^E(p_m;M)
  =
  \int b(m)\mathbf 1\{r_i^*(q,m;M,p_m,\ell_i)=E\}\,dF(q,m).
  \label{eq:p16m-project-demand}
\end{equation}
Study-related and total demand are
\begin{equation}
  \begin{aligned}
  D_m^{\mathrm{MAH}}(p_m;M)
  &=
  \int a_ix_i^*(M,p_m)\chi_i^E(p_m;M)\,dH(a,k,\ell),\\
  D_m(p_m;M)
  &=D_m^B(p_m)+D_m^{\mathrm{MAH}}(p_m;M).
  \end{aligned}
  \label{eq:p16m-demand}
\end{equation}
The demand expression includes both feedbacks: price changes expected value
and hence advancement, and it changes deterministic route selection.
Continuous heterogeneity makes aggregate demand continuous even though
individual indicators can jump.

The qualified-capacity market clears at the unique price
\begin{equation}
  D_m(p_m^*;M)=S_m(p_m^*).
  \label{eq:p16m-clearing}
\end{equation}
Supply is strictly increasing and demand weakly decreasing. Positive
background demand gives excess demand at zero, while entrusted and
background demand vanish and supply diverges at high prices. These conditions
give existence and uniqueness. If post-MAH study demand is positive at the
old price, \(p_m^*(1)>p_m^*(0)\): CMO scarcity attenuates the reform rather
than amplifying it through a direct price reduction.

For the equilibrium comparison, let \(p_m^h\equiv p_m^*(h)\), and define
\begin{align}
  \Delta\Omega_i^{\mathrm{dir}}
  &\equiv \Omega_i(1,p_m^0)-\Omega_i(0,p_m^0),
  &
  \Delta\Omega_i^{\mathrm{eq}}
  &\equiv \Omega_i(1,p_m^1)-\Omega_i(0,p_m^0),
  \label{eq:p16m-value-comparisons}\\
  \Delta x_i^{\mathrm{dir}}
  &\equiv x_i^*(1,p_m^0)-x_i^*(0,p_m^0),
  &
  \Delta x_i^{\mathrm{eq}}
  &\equiv x_i^*(1,p_m^1)-x_i^*(0,p_m^0).
  \label{eq:p16m-advancement-comparisons}
\end{align}

\begin{proposition}[CMO equilibrium and scarcity attenuation]
\label{prop:main-finance-cmo}
Under atomless financing and route heterogeneity, the existing integrability
conditions, weakly decreasing demand, and strictly increasing supply, the
scalar CMO market has a unique clearing price. Moreover,
\begin{equation}
  p_m^*(1)\geq p_m^*(0),
  \qquad
  0\leq\Delta\Omega_i^{\mathrm{eq}}
  \leq\Delta\Omega_i^{\mathrm{dir}},
  \qquad
  0\leq\Delta x_i^{\mathrm{eq}}
  \leq\Delta x_i^{\mathrm{dir}}.
  \label{eq:p16m-scarcity-attenuation}
\end{equation}
The price inequality is strict when post-MAH study demand is positive at the
old price. Financing capacity shifts neither CMO supply nor the background
demand schedule.
\end{proposition}

The baseline partial equilibrium contains exactly
\begin{equation}
  \left\{
  p_m^*,\;
  x_i^*,\;
  r_i^*(q,m),\;
  s_j^*
  \right\}_{i,j}.
  \label{eq:p16m-equilibrium}
\end{equation}
No entry, labor, capital, product-market aggregate, or welfare closure is
added.

\subsection{Observed outcomes and nesting}
\label{subsec:p16m-predictions}

\paragraph{Planning, organization, and observed outcomes.}
At equilibrium,
\[
  \Lambda_i^{\mathrm{plan}}=a_ix_i^*,
\]
while expected realized retained outcomes are
\begin{equation}
  Y_i^{\mathrm{ret}}
  =
  a_ix_i^*
  \int s(q)
  \left[
  \mathbf 1\{r_i^*=I\}
  +
  \mathbf 1\{r_i^*=E\}
  \right]dF(q,m).
  \label{eq:p16m-retained-outcome}
\end{equation}
Thus observed outcomes can change through advancement or route composition.
Holder--producer separation is observed only after route assignment and
cannot cause the earlier \(x_i\) choice.

\paragraph{Novelty composition.}
For \(g\in\{O,\mathrm{Inc}\}\), outcomes use the same \(x_i^*\) and the
mixture \(\rho_gF_g\). The baseline permits either class to have a zero or
larger response and imposes no original-versus-incremental ranking. Patent
applications are not a baseline endogenous outcome.

\paragraph{No-binding-finance limit.}
If
\[
  \ell_i\geq\max\{J_I(m,k_i),J_E(m)\}
\]
for every relevant developer--project pair, financing never binds,
\(\widetilde W_i^I=W_i^I\) and \(\widetilde W_i^E=W_i^E\). The model then
collapses to the previous no-financing-friction commercialization baseline.

\section{Empirical Predictions and Mapping Boundaries}

The deterministic model yields three conditional predictions. First, among
projects for which internal and entrusted production dominate transfer and
abandonment, developers with weaker internal manufacturing capability are
more likely to select retained entrusted production. This tendency is
strongest when pre-policy financing capacity lies between the entrusted and
internal up-front requirements. Second, MAH raises project advancement only
when the entrusted route is financeable and exceeds the best pre-reform
option; the response can be zero when financing is extremely scarce, when
internalization is already attractive and financeable, or when an outside
option dominates. Third, reform-induced demand for qualified manufacturing
services can raise \(p_m^*\), attenuating fixed-price gains in expected
project value and advancement. These predictions concern commercialization
organization and project advancement, not patent generation or scientific
productivity.

The empirical distinction is between ex ante advancement and downstream
observations. Product-level holder--manufacturer separation can proxy route
\(E\) only when holder identity, manufacturer identity, product identity, and
effective date are aligned. Approval or launch records are realized-product
outcomes, not direct observations of \(x_i\). Clinical milestones may proxy
project advancement only when their status and timing are measured consistently.
Patent applications are a separate patenting margin and are not a baseline
endogenous outcome. They may be studied as a separate empirical outcome, but
the present mechanism neither maps them to \(x_i\) nor assigns their reform
response a sign.

\begin{table}[htbp]
\centering
\footnotesize
\caption{Model Objects and Measurable Empirical Interfaces}
\label{tab:model_data_mapping}
\begin{tabularx}{\textwidth}{L{0.24\textwidth}Y Y}
\toprule
Model object & Candidate interface & Claim boundary \\
\midrule
Predetermined capability \(a_i\) & Pre-policy development history or other validated capability measure & Does not separately identify \(a_i\), and no policy shift in \(a_i\) is allowed. \\
Advancement \(x_i\) and \(a_ix_i\) & Consistently measured development-stage activity and planning-stage project events & \(x_i\) is not pure clinical effort, patent applications, or an approval count. \\
Predetermined financing capacity \(\ell_i\) & Validated pre-policy registered or paid-in capital, liquid assets, leverage or debt capacity, listed status, VC/PE or parent backing, or external-financing history & A proxy must predate reform exposure; it does not identify a finance market or imply that MAH changes financing supply. \\
Route \(r_i^*=E\) & Product-level holder differs from manufacturer, with time-consistent identifiers & The observed split is downstream evidence on retained organizational assignment, not the policy treatment itself. \\
Retained outcome \(Y_i^{\mathrm{ret}}\) & Approval or launch with developer/holder identity & Combines advancement, project mix, deterministic route assignment, and \(s(q)\); it does not identify these primitives separately. \\
CMO scarcity \(p_m^*\) & Service price, or consistently measured capacity, utilization, density, waiting-time, or participation proxy & A nonprice proxy does not identify the price or the supply and demand schedules separately. \\
Novelty class \(g\) & Versioned regulatory classification \(O\) or \(\mathrm{Inc}\) & The classifier creates no \(x_{ig}\) and imposes no response ranking. \\
Internal financing requirement \(J_I(m,k_i)\) & Validated up-front capital needed to build, acquire, validate, or upgrade compliant internal production & It is a liquidity threshold, not an extra route cost; no data source is asserted available before audit. \\
Entrusted financing requirement \(J_E(m)\) & Validated up-front transfer, validation, contracting, and holder-side preparation commitment & It is positive in the baseline and distinct from \(F_E,p_mb(m),\mu_E\); no data source is asserted available before audit. \\
\bottomrule
\end{tabularx}
\end{table}

The timing restriction can be summarized as
\begin{equation}
  M
  \longrightarrow
  \text{anticipated availability and value of }E
  \longrightarrow
  \Omega_i
  \longrightarrow
  x_i^*,
  \label{eq:empirical-anticipated-chain}
\end{equation}
followed by
\begin{equation}
  \begin{aligned}
  x_i^*
  &\longrightarrow
  \text{planning-stage projects}
  \longrightarrow r_i^* \\
  &\longrightarrow
  \text{observed holder--producer separation}
  \longrightarrow
  \text{realized products}.
  \end{aligned}
  \label{eq:empirical-realized-chain}
\end{equation}
Observed separation cannot be placed before or treated as a cause of the ex
ante \(x_i\) choice. The original and incremental classes use the same
\(x_i\) and the mixture \(\rho_gF_g\); either class may have a zero or larger
response.

These entries are interfaces for future validated data. No data set is
asserted here to be available, complete, or identifying. Before use, a source
must pass provenance, coverage, key-uniqueness, unit, timing, and
construct-validity checks. Entry, route-specific realization, and smooth
route choice remain undeveloped extensions. The primary financing prediction
is the triple interaction between MAH exposure, weak internal manufacturing
capability, and limited or intermediate pre-policy financing capacity, with
retained holder--manufacturer separation as the first outcome.

\section{Conclusion}

MAH changes expected project value by making retained entrusted manufacturing
available while preserving authorization-holder responsibility. The route is
most relevant when a developer can advance valuable projects but lacks
project-compatible internal manufacturing capability. Predetermined financing
capacity adds a feasibility margin: very low-capacity developers may be unable
to fund either retained route, intermediate-capacity developers may finance
entrusted production but not vertical integration, and high-capacity
developers recover internalization as an old option. MAH therefore changes
the organizational technology of commercialization; it does not relax the
financing constraint itself. Its equilibrium value also depends on
qualified-service scarcity: reform-induced entrusted demand can raise
\(p_m^*\) and attenuate the fixed-price gain. Project advancement rises only
for the financing-adjusted MAH-relevant set and need not rise for every
developer or novelty class. Observed holder--producer separation and realized
products occur after advancement and organizational assignment. Patent
applications and upstream scientific research remain outside the baseline,
and no patent response is imposed. The model makes no welfare claim.

\clearpage
\bibliographystyle{aer}
\bibliography{mah_route_indicator_friction_refs}

\end{document}
